\documentclass[11pt]{article}
\usepackage[utf8]{inputenc}
\usepackage[T1]{fontenc}
\usepackage{amsmath,amssymb,amsthm}
\usepackage{enumitem}
\usepackage[margin=1in]{geometry}

\newtheorem{definition}{Definition}
\newtheorem{theorem}{Theorem}
\newtheorem{remark}{Remark}

\newcommand{\I}{\mathcal{I}}
\newcommand{\E}{\mathcal{E}}
\newcommand{\MO}{M_\Omega}
\newcommand{\nullphi}{\varnothing_\phi}
\newcommand{\Omegaint}{\Omega_{\mathrm{int}}}
\newcommand{\Ext}{\mathrm{Ext}}

\title{CCR Null-Regime Conditional Closure v31}
\author{QICN Internal Hardening Artifact}
\date{2026-05-29}

\begin{document}
\maketitle

\noindent\textbf{Governance boundary.}
This document is an internal mathematical hardening artifact. It does not certify external support, consciousness, phenomenality, identity transfer, bridge-burden closure, or human mathematical review.

\begin{definition}[Typed perturbation boundary]
For a system $S$ with identity object $\I(S)$, let $\Omegaint(S)$ denote finite-energy perturbations whose domain and codomain remain inside the admissible identity structure of $S$. Let $\Ext(S)$ denote extension-side witnesses that compare $\I(S)$ with an independently specified object $\tilde{\I}$ not introduced as an element of $\Omegaint(S)$.
\end{definition}

\begin{definition}[Undefined-assignment marker and separated external witness]
The symbol $\bot$ (``bottom'') denotes absence of an admissible phenomenological assignment. It is disjoint from $\E$, so $\bot\notin\E$. This is distinct from the null regime $\nullphi\in\E$, which is the bottom element of the phenomenological poset in Paper~3. A separated external witness for a phenomenological assignment $\Phi$ is a tuple $(\tilde{\I},\epsilon,C)$ such that:
\begin{enumerate}[leftmargin=*]
\item $\tilde{\I}\in\Ext(S)$ and $d_{\I}(\I(S),\tilde{\I})=\epsilon>0$;
\item $C>0$ is a lower Lipschitz constant on the declared two-point witness domain; and
\item $\Phi$ is defined and admissible on both $\I(S)$ and $\tilde{\I}$, so $\Phi(\I(S)),\Phi(\tilde{\I})\in\E$.
\end{enumerate}
\end{definition}

\begin{theorem}[Witness-relative undefined-assignment exclusion]
Let $S$ be a CCR channel with $\MO=+\infty$. Let $\Phi$ be a declared compatibility operator satisfying the lower Lipschitz condition on the separated witness domain. Suppose a separated external witness $(\tilde{\I},\epsilon,C)$ exists for $\Phi$. Then $\Phi$ cannot be undefined on both witness endpoints:
\[
  \neg(\Phi(\I(S))=\bot \wedge \Phi(\tilde{\I})=\bot).
\]
\end{theorem}

\begin{proof}
Suppose for contradiction that both $\Phi(\I(S))=\bot$ and $\Phi(\tilde{\I})=\bot$.
By the definition of $\bot$, neither endpoint has an admissible $\E$-valued phenomenological assignment. This contradicts the separated-witness hypothesis, which requires $\Phi$ to be defined and admissible on both endpoints, so that $\Phi(\I(S)),\Phi(\tilde{\I})\in\E$ and the lower Lipschitz condition is meaningful there.

The case where both endpoints are assigned the null regime $\nullphi\in\E$ is a different typed statement: it is governed by the Paper~3 instability theorem and uses the metric on $\E$. The present theorem concerns only the absence marker $\bot\notin\E$.
\end{proof}

\begin{remark}[Non-claim]
CCR alone does not produce a witness, does not prove that a compatibility operator exists, and does not imply consciousness or experience. The theorem is conditional on a typed external witness and a checked lower Lipschitz bound. In particular, $\bot$ is not a point in $\E$; it marks absence of assignment. The null regime $\nullphi\in\E$ is a separate concept: it is the bottom element of the phenomenological poset, not the absence of a value.
\end{remark}

\end{document}
